% GENERATED FILE. Edit canonical lessons, then run npm run generate:books.
% canonical-source-hash: fnv1a64:5da49c769ea37546
% canonical-lessons: BN-C439-sukh, BN-C439-koshto, BN-C439-tripti, BN-C439-upobhog-kora, BN-C439-mete-otha

\chapter{Happiness and Hardship}
\label{ch:bn-happiness-and-hardship}
\hlchaptermodality{439}

\begin{chapteropening}
\textbf{By the end of this chapter:} \emph{I can talk about happiness, hardship and enjoying something.}

Complete the last lesson of chapter 439: I can talk about happiness, hardship and enjoying something.
\end{chapteropening}

\section[\texorpdfstring{\bn{সুখ} (sukh)}{সুখ (sukh)}]{\bn{সুখ} (sukh) — happiness; comfort}

\label{lesson:BN-C439-sukh}



\begin{quote}
Before the new one: say the Bengali for to feel, then the Bengali for to become aware of.
\end{quote}

\subsection*{You'll want to know: \bn{সুখ}}
\textbf{\bn{সুখ}} — \emph{sukh} — \textquotedblleft{}happiness; comfort\textquotedblright{}.

\begin{quote}
\bn{আমার} \bn{দাদু} \bn{গ্রামে} \bn{সুখে} \bn{আছেন।}
\end{quote}

\begin{quote}
\emph{āmār dādu grāme sukhe āchhen.}
\end{quote}

\begin{quote}
My grandfather lives happily in the village.
\end{quote}

\bn{সুখে} \bn{থাকা}, to live happily.

\begin{cousinweb}
\textbf{\bn{সুখ}} (\emph{sukh}) is a tatsama, Sanskrit \emph{sukha}, \textquotedblleft{}pleasant, easy\textquotedblright{} (CDIAL 13451), and so comfort and happiness; \bn{সুখী} (taught), contented, is the adjective made from it.
\end{cousinweb}

\begin{grammarlens}[title={What you've built}]
One more word for describing what you lived and felt.
\end{grammarlens}

\subsection*{Your turn}
\begin{itemize}
\raggedright
  \item \emph{Say it:} \emph{sukh}
  \item \emph{Say it:} \emph{sukh}, once more
  \item \emph{Say it:} \emph{āmār dādu grāme sukhe āchhen.} — My grandfather lives happily in the village.
  \item \emph{Recall:} say the Bengali for to feel, then the Bengali for to become aware of, then say \emph{sukh} again
\end{itemize}

\begin{tcolorbox}[breakable,colback=teal!4,colframe=teal!35!black,title={Before you move on}]
What does \bn{সুখ} mean? (\textquotedblleft{}Happiness; comfort\textquotedblright{}.) Say it once more.
\end{tcolorbox}
\section[\texorpdfstring{\bn{কষ্ট} (kôshṭô)}{কষ্ট (kôshṭô)}]{\bn{কষ্ট} (kôshṭô) — suffering; hardship}

\label{lesson:BN-C439-koshto}



\begin{quote}
Before the new one: say the Bengali for to become aware of, then the Bengali for happiness.
\end{quote}

\subsection*{You'll want to know: \bn{কষ্ট}}
\textbf{\bn{কষ্ট}} — \emph{kôshṭô} — \textquotedblleft{}suffering; hardship\textquotedblright{}.

\begin{quote}
\bn{গরমে} \bn{বাসে} \bn{যেতে} \bn{আমার} \bn{কষ্ট} \bn{হয়।}
\end{quote}

\begin{quote}
\emph{gôrôme bāse jete āmār kôshṭô hôy.}
\end{quote}

\begin{quote}
Travelling by bus in the heat is hard for me.
\end{quote}

\bn{কষ্ট} \bn{হওয়া}, to be hard for someone: \bn{আমার} \bn{কষ্ট} \bn{হয়}. \bn{কষ্ট} \bn{পাওয়া}, to suffer, to be hurt.

\begin{cousinweb}
\textbf{\bn{কষ্ট}} (\emph{kôshṭô}) is a tatsama, Sanskrit \emph{kaṣṭa}, \textquotedblleft{}bad, ill, painful\textquotedblright{} (Monier-Williams); Bengali uses it as a noun, the hardship itself.
\end{cousinweb}

\begin{grammarlens}[title={What you've built}]
One more word for describing what you lived and felt.
\end{grammarlens}

\subsection*{Your turn}
\begin{itemize}
\raggedright
  \item \emph{Say it:} \emph{kôshṭô}
  \item \emph{Say it:} \emph{kôshṭô}, once more
  \item \emph{Say it:} \emph{gôrôme bāse jete āmār kôshṭô hôy.} — Travelling by bus in the heat is hard for me.
  \item \emph{Recall:} say the Bengali for to become aware of, then the Bengali for happiness, then say \emph{kôshṭô} again
\end{itemize}

\begin{tcolorbox}[breakable,colback=teal!4,colframe=teal!35!black,title={Before you move on}]
What does \bn{কষ্ট} mean? (\textquotedblleft{}Suffering; hardship\textquotedblright{}.) Say it once more.
\end{tcolorbox}
\section[\texorpdfstring{\bn{তৃপ্তি} (tripti)}{তৃপ্তি (tripti)}]{\bn{তৃপ্তি} (tripti) — satisfaction; contentment}

\label{lesson:BN-C439-tripti}



\begin{quote}
Before the new one: say the Bengali for happiness, then the Bengali for suffering.
\end{quote}

\subsection*{You'll want to know: \bn{তৃপ্তি}}
\textbf{\bn{তৃপ্তি}} — \emph{tripti} — \textquotedblleft{}satisfaction; contentment\textquotedblright{}.

\begin{quote}
\bn{দিদিমার} \bn{রান্না} \bn{খেয়ে} \bn{আমরা} \bn{তৃপ্তি} \bn{পেলাম।}
\end{quote}

\begin{quote}
\emph{didimār rānnā kheye āmrā tripti pelām.}
\end{quote}

\begin{quote}
We ate my grandmother's cooking with real satisfaction.
\end{quote}

\bn{তৃপ্তি} \bn{পাওয়া}, to find satisfaction; \bn{তৃপ্তি} \bn{করে} \bn{খাওয়া}, to eat with relish.

\begin{cousinweb}
\textbf{\bn{তৃপ্তি}} (\emph{tripti}) is a tatsama, Sanskrit \emph{tṛpti}, \textquotedblleft{}satisfaction, contentment\textquotedblright{} (Monier-Williams): the feeling of having had enough.
\end{cousinweb}

\begin{grammarlens}[title={What you've built}]
One more word for describing what you lived and felt.
\end{grammarlens}

\subsection*{Your turn}
\begin{itemize}
\raggedright
  \item \emph{Say it:} \emph{tripti}
  \item \emph{Say it:} \emph{tripti}, once more
  \item \emph{Say it:} \emph{didimār rānnā kheye āmrā tripti pelām.} — We ate my grandmother's cooking with real satisfaction.
  \item \emph{Recall:} say the Bengali for happiness, then the Bengali for suffering, then say \emph{tripti} again
\end{itemize}

\begin{tcolorbox}[breakable,colback=teal!4,colframe=teal!35!black,title={Before you move on}]
What does \bn{তৃপ্তি} mean? (\textquotedblleft{}Satisfaction; contentment\textquotedblright{}.) Say it once more.
\end{tcolorbox}
\section[\texorpdfstring{\bn{উপভোগ} \bn{করা} (upôbhog kôrā)}{উপভোগ করা (upôbhog kôrā)}]{\bn{উপভোগ} \bn{করা} (upôbhog kôrā) — to enjoy}

\label{lesson:BN-C439-upobhog-kora}



\begin{quote}
Before the new one: say the Bengali for suffering, then the Bengali for satisfaction.
\end{quote}

\subsection*{You'll want to know: \bn{উপভোগ} \bn{করা}}
\textbf{\bn{উপভোগ} \bn{করা}} — \emph{upôbhog kôrā} — \textquotedblleft{}to enjoy\textquotedblright{}.

\begin{quote}
\bn{আমরা} \bn{ছুটিটা} \bn{পুরো} \bn{উপভোগ} \bn{করলাম।}
\end{quote}

\begin{quote}
\emph{āmrā chhuṭiṭā puro upôbhog kôrlām.}
\end{quote}

\begin{quote}
We enjoyed the holiday to the full.
\end{quote}

Past (he/she …): \bn{তিনি} \bn{গানটা} \bn{উপভোগ} \bn{করলেন}, \emph{tini gānṭā upôbhog kôrlen}.

\begin{cousinweb}
\textbf{\bn{উপভোগ} \bn{করা}} (\emph{upôbhog kôrā}) is Sanskrit \emph{upabhoga}, \textquotedblleft{}enjoyment, using\textquotedblright{} (Monier-Williams), with \emph{kôrā}, \textquotedblleft{}to do\textquotedblright{} (CDIAL 2814).
\end{cousinweb}

\begin{grammarlens}[title={What you've built}]
One more word for describing what you lived and felt.
\end{grammarlens}

\subsection*{Your turn}
\begin{itemize}
\raggedright
  \item \emph{Say it:} \emph{upôbhog kôrā}
  \item \emph{Say it:} \emph{upôbhog kôrā}, once more
  \item \emph{Say it:} \emph{āmrā chhuṭiṭā puro upôbhog kôrlām.} — We enjoyed the holiday to the full.
  \item \emph{Recall:} say the Bengali for suffering, then the Bengali for satisfaction, then say \emph{upôbhog kôrā} again
\end{itemize}

\begin{tcolorbox}[breakable,colback=teal!4,colframe=teal!35!black,title={Before you move on}]
What does \bn{উপভোগ} \bn{করা} mean? (\textquotedblleft{}To enjoy\textquotedblright{}.) Say it once more.
\end{tcolorbox}
\section[\texorpdfstring{\bn{মেতে} \bn{ওঠা} (mete oṭhā)}{মেতে ওঠা (mete oṭhā)}]{\bn{মেতে} \bn{ওঠা} (mete oṭhā) — to get carried away; to be swept up in excitement}

\label{lesson:BN-C439-mete-otha}



\begin{quote}
Before the new one: say the Bengali for satisfaction, then the Bengali for to enjoy.
\end{quote}

\subsection*{You'll want to know: \bn{মেতে} \bn{ওঠা}}
\textbf{\bn{মেতে} \bn{ওঠা}} — \emph{mete oṭhā} — \textquotedblleft{}to get carried away; to be swept up in excitement\textquotedblright{}.

\begin{quote}
\bn{খেলা} \bn{দেখে} \bn{ছেলেরা} \bn{আনন্দে} \bn{মেতে} \bn{উঠল।}
\end{quote}

\begin{quote}
\emph{khælā dekhe chhelerā ānônde mete uṭhlô.}
\end{quote}

\begin{quote}
Watching the game, the boys got carried away with joy.
\end{quote}

Past (he/she …): \bn{তিনি} \bn{গল্পে} \bn{মেতে} \bn{উঠলেন}, \emph{tini gôlpe mete uṭhlen}.

\begin{cousinweb}
\textbf{\bn{মেতে} \bn{ওঠা}} (\emph{mete oṭhā}) is \emph{mete}, the -\bn{এ} form of \emph{mātā}, \textquotedblleft{}to be wild with excitement\textquotedblright{}, which belongs with the Bengali \emph{māt}, \textquotedblleft{}delight\textquotedblright{}, that CDIAL lists under a reconstructed form (CDIAL 9753), with \emph{oṭhā}, \textquotedblleft{}to rise\textquotedblright{} (CDIAL 1793).
\end{cousinweb}

\begin{grammarlens}[title={What you've built}]
One more word for describing what you lived and felt.
\end{grammarlens}

\subsection*{Your turn}
\begin{itemize}
\raggedright
  \item \emph{Say it:} \emph{mete oṭhā}
  \item \emph{Say it:} \emph{mete oṭhā}, once more
  \item \emph{Say it:} \emph{khælā dekhe chhelerā ānônde mete uṭhlô.} — Watching the game, the boys got carried away with joy.
  \item \emph{Recall:} say the Bengali for satisfaction, then the Bengali for to enjoy, then say \emph{mete oṭhā} again
\end{itemize}

\begin{tcolorbox}[breakable,colback=teal!4,colframe=teal!35!black,title={Before you move on}]
What does \bn{মেতে} \bn{ওঠা} mean? (\textquotedblleft{}To get carried away; to be swept up in excitement\textquotedblright{}.) Say it once more.
\end{tcolorbox}
